% sec13_1.tex — Section 13.1: Introduction
\section{Introduction}
\label{sec:13.1}

We have introduced some techniques to solve linear partial differential equations.
For problems with a simple geometry, the method of separation of variables motivates using Fourier series, its various generalizations, or variants of the Fourier transform.  Of most importance is the type of boundary condition, including whether the domain is finite, infinite or semi-infinite.  In some problems a Green's function can be utilized, while for the one-dimensional wave equation the method of characteristics exists.  Whether or not any of these methods may be appropriate, numerical methods (introduced in Chapter~6) are often most efficient.

Another technique, to be elaborated on in this chapter, relies on the use of Laplace transforms.  \emph{Most problems in partial differential equations that can be analyzed by Laplace transforms also can be analyzed by one of our earlier techniques, and substantially equivalent answers can be obtained.}  The use of Laplace transforms is advocated by those who feel more comfortable with them than with our other methods.  Instead of taking sides, we will present the elementary aspects of Laplace transforms in order to enable the reader to become somewhat familiar with them.  However, whole books\footnote{For example, Churchill~[1972].} have been written concerning their use in partial differential equations.  Consequently, in this chapter we only briefly discuss Laplace transforms and describe their application to partial differential equations with only a few examples.
